\section{Data, models and how they were compared}\label{sec:methods}

\paragraph{Data}
The analysis uses the 12,937 trials of monkey Sirius that Vera selected: 42 sessions from 30 March to 31 July 2023, both hands tested in every session, about 80 trials per hand and condition in each session.
A vibration of 6--50\um\ is applied to one hand and the monkey reports whether it is stronger or weaker than the trained boundary of 28\um.
On half of the trials (bimanual) the other hand receives a distractor vibration at the same time.
The distractor amplitude is fixed within a session and was raised over the study, from 6\um\ in the first eight sessions to 32\um\ in most of the later ones (Fig.~\ref{fig:time}, top).

\paragraph{Models}
All models share one psychometric curve with ``high'' and ``low'' lapse rates (Appendix~\ref{app:models}).
They differ in what is allowed to change between hands, conditions and sessions.
Vera's Model~B fits unimanual and bimanual trials of each hand as unrelated data sets.
Her Model~C fits one curve per hand and shifts it on bimanual trials by an amount proportional to the distractor amplitude, with the same lapse rates in both conditions.
Model~D, recommended in the previous report, is C with three additions: bimanual trials get their own lapse rates (L), these change with the distractor amplitude (K), and a distractor of any amplitude can shift the curve (presence, P).
This report adds simpler versions of D, which keep parameters fixed across sessions, and extensions of D that ask whether the monkey's criterion moved between sessions in a systematic way. The recommended Model~E combines one of each.
Table~\ref{tab:comparison} lists all \nModels\ models.

\paragraph{Sampling}
Every model was refit with Vera's setting: 4 chains of 5000 draws after 1000 tuning steps, target acceptance 0.9.
Model~C and its variants were sampled with nutpie, a faster implementation of the same NUTS algorithm; A and the B family with PyMC's own sampler, which handles B's parameterization better.
The model code is the same in both cases, and every model's likelihood was checked against an independent implementation.
All fits are clean (Appendix~\ref{app:diag}): \totalDiv\ divergences in total, $\hat R\le\maxRhatAll$, and at least \minEssAll\ effective draws for every parameter.

\paragraph{Two ways to score a model}
Both scores estimate the expected log predictive density (elpd) of data the model has not seen; higher is better, and differences are in log-likelihood units summed over trials.
\begin{itemize}
\item \textbf{Trial-level LOO} is the criterion in Vera's reports (PSIS-LOO). It predicts each trial from all the others, including the other trials of the same session, so it asks how well the model describes the sessions it was fit to.
Trials from one session share parameters, so the usual standard error, which treats trials as independent, is unreliable (for the larger differences here it is too small); the standard error of every difference is computed from the 42 per-session sums instead.
\item \textbf{Leave-one-session-out (LOSO)} predicts every trial of a session from the other 41 sessions, with that session's own parameters drawn from the population distribution.
It asks how well the model would predict a new session, which is what claims about the population-level effects (the shift per\um\ of distractor, the presence offset) rest on.
It is computed without refitting, by importance sampling over the held-out session's parameters followed by PSIS over the posterior draws (1000 draws thinned from the 20,000).
In checks with exact refits without the held-out session, the two agreed to within 0.2 elpd, including sessions with Pareto $k$ up to 1.03 (Appendix~\ref{app:diag}).
\end{itemize}
A difference is treated as clear when it exceeds about twice its standard error.
